\item En el arreglo que se muestra en la figura P17.14, un objeto de masa $m = 5.00\text{ kg}$ cuelga de una cuerda alrededor de una polea ligera. La longitud del cable entre la polea y el punto P es $L = 2.00\text{ m}$.
\begin{enumerate}
    \item Cuando el vibrador se ajusta a una frecuencia de $150\text{ Hz}$, se forma una onda estacionaria con seis bucles. ¿Cuál debe ser la densidad de masa lineal de la cuerda?
    \item ¿Cuántos bucles (si existen) se producirán si $m$ se cambia a $45.0\text{ kg}$?
    \item ¿Cuántos bucles (si existen) se producirán si $m$ se cambia a $10.0\text{ kg}$?
\end{enumerate}